\scanpage{116}
\mspage{110}
% archive sheet number: 115

Sous ces conditions, \deleted{$\unsure{B|\Sigma}$} l'image $\Sigma_e$ de cette application est le plus petit sous-ensemble $\Sigma'$ de $\Sigma$ satisfaisant les conditions de (i).

\uline{Cor 1} \unsure{Mêmes|Memes} \unsure{données|donnees}, mais on ne suppose pas que $e\overset{\mathrm{def}}{=}\operatorname{Sup}_{i\in I}e_i$ soit égal à $1$, et on introduit $e_0=\unsure{\complement|C}e$. Pour que les conditions du lemme soient satisfaites pour $\{\{e_i\},e_0\}$ (indexé pr $I\amalg\{0\}$) il f et s que \deleted{l'on ait,} pour tte partie $I'$ de $I$, posant $I''=I\setminus I'$, (de sorte que $e=e_{I'}\vee e_{I''}$, \deleted{on ait} $e_0=\complement e$), \deleted{on ait} \uline{on} ait
\[
\complement e_{I'}=e_{I''}\vee e_0
\]
ou ce qui revient au $=$ (\unsure{passant} aux $\complement$)
\[
e_{I'}=e\wedge\complement e_{I''}.
\]

\uline{Def} On dit alors que la \deleted{système}\inserted{famille} $(e_i)_{i\in I}$ d'éléments mutuellement disjoints est ``\uline{admissible}''. $\Sigma_{e_{\unsure{\ast|i|\bullet}}}=\Sigma_{(e_i)}$, alg. de Boole \unsure{engendrée}, \unsure{is.|iso.} à $\unsure{\mathscr{P}|\mathcal{P}|P}(I\amalg\{e_0\})$.
